\documentclass[11pt]{article}
\usepackage{amsmath}
\usepackage{booktabs}
\usepackage{caption}
\usepackage[margin=1in]{geometry}
\usepackage{graphicx}
\usepackage[sc]{mathpazo}
\usepackage{setspace}
\usepackage[table]{xcolor}
\usepackage[colorlinks,linkcolor={black!50!black},citecolor={blue!50!black},urlcolor={blue!80!black}]{hyperref}

\title{Task 4: Foreign Non-Profits and Tax-Haven Owners \\
\large Descriptive Profiles, Outliers, and Case Studies}
\author{Disconnected Landlords Project}
\date{\today}

\begin{document}
\maketitle

\begin{abstract}
\noindent Descriptive deep-dive into two ownership groups flagged in the PI
feedback: foreign-incorporated non-profit proprietors and tax-haven
incorporated proprietors of privately rented housing in England and Wales.
We characterise the universe (counts by local authority, jurisdiction of
incorporation, and region), identify the largest proprietors and outlier
properties, and develop a shortlist of case studies from public-register
data. All named entities are registered companies from the public CCOD/OCOD
registers; no individual beneficial owners are identified.
\end{abstract}

\section{Data and definitions}

Groups are defined exactly as the corresponding regression treatments
(\texttt{treatment\_definitions.R}):
\begin{itemize}
  \item \textbf{Foreign non-profit} (\texttt{frnp}): proprietorship category
    Non-Profit/Community Organisations, incorporated outside the United Kingdom.
  \item \textbf{Tax-haven non-profit} (\texttt{thnp}): non-profit category,
    incorporated in an IMF-list tax-haven jurisdiction.
  \item \textbf{Tax haven, all} (\texttt{th}): any proprietor incorporated in a
    tax-haven jurisdiction, with sub-classes British (Crown dependencies and
    overseas territories), European, Caribbean, and other.
\end{itemize}
The reference group is the standard control: privately rented properties with
unknown ownership (not in CCOD/OCOD). Properties are deduplicated by UPRN;
title-to-UPRN expansion means title counts and property counts are reported
separately. Produced by \texttt{scripts/08b\_offshore\_nonprofit\_profiles.R},
outputs in \texttt{output/offshore\_profiles/}.

\section{Universe}

Across England the register links 181{,}952 domestic properties to a
tax-haven-incorporated or foreign-incorporated proprietor: 174{,}947 to a
tax-haven company (group \texttt{th}), of which 10{,}777 sit under a foreign
non-profit (\texttt{frnp}) and 3{,}772 under a tax-haven non-profit
(\texttt{thnp}). The non-profit groups overlap the tax-haven group by
construction (a Jersey housing society is both foreign and a tax haven), so the
three counts are not additive.

The universe is overwhelmingly held through the British Crown Dependencies and
Overseas Territories: 167{,}996 properties (92 percent) are incorporated in a
``british'' haven, led by Guernsey (76{,}691), Jersey (56{,}714), the British
Virgin Islands (19{,}746) and the Isle of Man (11{,}790). This concentration is
an artefact of who uses these jurisdictions: the Channel Islands are the home of
the large ground-rent freehold vehicles (see below), while the classic
wealth-holding havens supply a much smaller but distinctive tail.

A clear efficiency gradient runs across jurisdictions. The Channel Islands
freehold stock is efficient and new (Guernsey mean EPC 78, 11 percent below band
C; Jersey 77, 14 percent), because these owners hold only the freehold reversion
of modern leasehold flats. The opaque wealth havens hold older, far less
efficient prime and period stock: the British Virgin Islands (mean 72, 28
percent below C), Isle of Man (72, 30 percent), Gibraltar (67, 49 percent),
Panama (66, 48 percent) and the Seychelles (68, 43 percent). Aggregated by
class, ``british''-incorporated properties are 16 percent below band C against
40 percent for ``caribbean'' and 37 percent for ``european'' incorporations. The
foreign non-profit tail is dominated by Luxembourg (5{,}384 properties, mean EPC
78), reflecting institutional student and build-to-rent vehicles rather than
wealth holding.

\begin{table}[htbp]
\centering
\caption{Tax-Haven Owned Rental Properties by Jurisdiction of Incorporation}
\label{tab:offshore_universe_jurisdiction}
\small
\begin{tabular}{lrrrrr}
\toprule
Jurisdiction & Properties & Titles & Proprietors & Mean EPC & Share $<$C \\
\midrule
GUERNSEY & 76,691 & 5,298 & 784 & 78.2 & 10.8\% \\
JERSEY & 56,714 & 5,036 & 1,310 & 77.2 & 14.5\% \\
BRITISH VIRGIN ISLANDS & 19,746 & 6,811 & 3,825 & 72.4 & 28.1\% \\
ISLE OF MAN & 11,790 & 3,249 & 1,192 & 71.7 & 30.3\% \\
GIBRALTAR & 1,999 & 737 & 278 & 66.5 & 48.9\% \\
HONG KONG & 983 & 566 & 313 & 74.3 & 20.5\% \\
CYPRUS & 872 & 436 & 243 & 72.8 & 30.4\% \\
PANAMA & 735 & 375 & 278 & 66.5 & 48.2\% \\
BAHAMAS & 707 & 300 & 189 & 69.9 & 36.5\% \\
SEYCHELLES & 633 & 446 & 314 & 67.9 & 43.0\% \\
CAYMAN ISLANDS & 489 & 206 & 124 & 72.9 & 27.8\% \\
BERMUDA & 378 & 76 & 42 & 58.9 & 90.7\% \\
BELIZE & 295 & 158 & 86 & 73.5 & 27.1\% \\
MAURITIUS & 272 & 220 & 149 & 68.5 & 37.9\% \\
LIECHTENSTEIN & 239 & 136 & 75 & 63.0 & 54.8\% \\
\bottomrule
\end{tabular}
\begin{minipage}{0.9\textwidth}
\vspace{4pt}
\footnotesize
\textit{Notes:} Properties in the tax-haven ownership group (IMF tax-haven
jurisdiction lists applied to CCOD/OCOD country of incorporation), deduplicated
by UPRN. Top 15 jurisdictions by property count. Share $<$C = share of properties
below EPC band C (SAP $<$ 69). Source: \texttt{universe\_by\_jurisdiction.csv} (script 08b).
\end{minipage}
\end{table}


\section{Largest proprietors}

\begin{table}[htbp]
\centering
\caption{Largest Tax-Haven and Foreign Non-Profit Proprietors (Public Register)}
\label{tab:offshore_top_proprietors}
\scriptsize
\begin{tabular}{llrrrr}
\toprule
Group & Proprietor (registered company) & Properties & LAs & Mean EPC & Share $<$C \\
\midrule
Tax haven (all) & ABACUS LAND 4 LIMITED (GUERNSEY) & 18,870 & 240 & 76.2 & 12.9\% \\
Tax haven (all) & ADRIATIC LAND 3 LIMITED (GUERNSEY) & 14,490 & 204 & 78.8 & 7.9\% \\
Tax haven (all) & ABACUS LAND 1 (HOLDCO 1) LIMITED (JERSEY) & 12,613 & 180 & 76.8 & 11.3\% \\
Tax haven (all) & PENNINE TRUSTEES LIMITED (JERSEY) & 11,216 & 48 & 79.2 & 6.4\% \\
Tax haven (all) & ADRIATIC LAND 11 LIMITED (GUERNSEY) & 8,544 & 152 & 78.7 & 9.3\% \\
Tax haven (all) & ADRIATIC LAND 5 LIMITED (GUERNSEY) & 6,625 & 112 & 78.9 & 12.3\% \\
Tax haven (all) & ADRIATIC LAND 12 LIMITED (GUERNSEY) & 3,624 & 54 & 82.2 & 3.6\% \\
Tax haven (all) & ADRIATIC LAND 9 LIMITED (GUERNSEY) & 3,189 & 72 & 81.6 & 3.5\% \\
Tax haven (all) & ADRIATIC LAND 10 LIMITED (GUERNSEY) & 3,038 & 53 & 81.9 & 8.3\% \\
Tax haven (all) & ADRIATIC LAND 7 LIMITED (GUERNSEY) & 2,825 & 78 & 80.9 & 5.0\% \\
Foreign non-profit & THE DOLPHIN SQUARE ESTATE S.A.R.L. (LUXEMBOURG) & 1,054 & 1 & 70.3 & 20.2\% \\
Foreign non-profit & BALTIMORE WHARF SLP (JERSEY) & 937 & 1 & 82.0 & 0.4\% \\
Foreign non-profit & FREEHOLD REVERSIONS PARTNERSHIP INCORPORATED  (GUERNSEY) & 503 & 15 & 79.2 & 3.9\% \\
Foreign non-profit & WHITEWOOD (RING) UK S.A.R.L. (LUXEMBOURG) & 421 & 1 & 85.7 & 0.0\% \\
Foreign non-profit & GREENGATE S.A R.L. (LUXEMBOURG) & 351 & 1 & 82.6 & 0.3\% \\
\bottomrule
\end{tabular}
\begin{minipage}{0.95\textwidth}
\vspace{4pt}
\footnotesize
\textit{Notes:} Proprietor names are registered company names from the public
CCOD/OCOD registers (company-level data only; no individual beneficial owners).
Properties = distinct UPRNs held under the (normalised) proprietor name.
Source: \texttt{top\_proprietors.csv} (script 08b).
\end{minipage}
\end{table}


Ownership is extraordinarily concentrated at the top, and the largest holders
are all Channel Islands ground-rent freehold vehicles. The ten largest
proprietors are led by Abacus Land~4 (Guernsey, 18{,}870 properties), Adriatic
Land~3 (Guernsey, 14{,}490), Abacus Land~1 Holdco~1 (Jersey, 12{,}613) and
Pennine Trustees (Jersey, 11{,}216), followed by a string of further Adriatic
Land vehicles (numbers 11, 5, 12 and 9). These are freehold-reversion investment
companies rather than occupier-landlords: they hold the ground rents beneath
modern leasehold flats, which is why their portfolios are large, geographically
dispersed (each spans hundreds of local authorities) and comparatively efficient
(mean EPC 76 to 82, only 4 to 13 percent below band C). This is the opposite
profile to the small opaque-haven holders, which own single prime or period
units with poor efficiency. Portfolio size and low efficiency are therefore
negatively, not positively, correlated in this universe: the biggest offshore
landlords are ground-rent financiers, not the worst-performing owners.

\section{Outliers}

The outlier table reports, per group, the five worst-EPC, largest-floor-area,
highest-price, largest-portfolio-member and most-unusual-combination properties
(75 rows). Implausible values (floor area under 10 or over 1{,}000 sqm, price
under 1{,}000) are retained but flagged \texttt{data\_quality\_suspect} rather
than presented as headline cases; the worst-EPC extremes in particular are
dominated by nominal-price freehold-reversion transfers (a ``price'' of a few
thousand pounds for a ground rent) and should not be read as market sales. The
substantive outliers, super-prime central London units held through wealth
havens with band F or G efficiency, are carried forward as the case studies
below.

\section{Case studies}

% One subsection per researched case; drafted from case_studies.csv shortlist
% plus public Companies House filings and press coverage.
% PRIVACY RULE: company-level public data only; no individual beneficial-owner
% names in any published output.

Each proprietor below is a registered company; facts are drawn from the UK
Register of Overseas Entities (ROE) and Companies House, and from reputable press
or public planning records where cited. Consistent with the project's disclosure
rule, only company-level and public-register information is reported: no
individual director, shareholder or beneficial-owner names appear. Register
numbers are the ROE ``OE'' identifiers on the Companies House service. Where a
source supports only a corporate-family link rather than the exact title-holding
vehicle, that is said plainly.

\subsection{Adriatic Land 3 (Guernsey): ground-rent freeholds}

Adriatic Land 3 Limited (ROE \texttt{OE018380}, incorporated in Guernsey,
registered on the ROE on 26 January 2023) is one of the Channel Islands vehicles
behind the largest offshore freehold holding in England and Wales. Our data
attributes to it roughly 14{,}490 properties across 321 local authorities, with a
good mean efficiency (EPC 79) characteristic of a freehold-reversion owner of
modern leasehold flats rather than an occupier-landlord. Reporting by Global
Witness places the Adriatic Land and Abacus Land companies together at 2{,}926
freeholds, describes them as the country's largest offshore corporate freeholders,
and links the group to the Long Harbour ground-rent investment
fund.\footnote{\url{https://www.globalwitness.org/en/campaigns/corruption-and-money-laundering/anonymous-company-owners/legalised-extortion/}}
In March 2022 Adriatic Land 3 Limited gave voluntary undertakings to the
Competition and Markets Authority to remove doubling ground-rent clauses at a
north London development, freeze ground rents at their commencement level, and
meet leaseholders'
costs.\footnote{\url{https://www.homegroundonline.com/news/2022/03/voluntary-undertakings-by-adriatic-land-3-limited-with-the-cma/}}
A sibling entity, Abacus Land 4, illustrates the RPI-indexed model, its ground
rent rising from \pounds{}350 to \pounds{}711 at first
review.\footnote{\url{https://www.leaseholdknowledge.com/abacus-land-sees-its-ground-rents-more-than-double-to-711-thanks-to-rpi/}}
(The Guernsey ROE entity is distinct from the separately UK-incorporated
``Adriatic Land 3 (GR1) Limited''; the specific Edgbaston freehold in our data is
not among the published CMA cases.)

\subsection{Vevil Properties (Jersey): the Curzon Mayfair estate}

Vevil Properties Limited (ROE \texttt{OE013879}, Jersey, registered 17 January
2023) belongs to the Vevil corporate family, reported to sit under Vevil
International Ltd in the British Virgin Islands (where beneficial ownership is not
publicly filed) with Jersey holding and operating
subsidiaries.\footnote{\url{https://londoninbits.substack.com/p/the-curse-of-the-curzon}}
A 2016 marketing brochure listed the family's central London portfolio as
including 38 Curzon Street, the site of the Curzon Mayfair cinema, alongside 160
Piccadilly (The Wolseley) and 27 Poultry (The Ned). The Curzon Mayfair (37--38
Curzon Street) is a Grade~II listed 1960s building described by Historic England
as the finest surviving post-war cinema, and it survived a 2016 closure threat
after a public campaign. In our data the associated title carries a band-G
efficiency (EPC 3), the extreme low end of the prime-central tail. The precise
mapping of ``Vevil Properties Limited'' to the Curzon title, as opposed to a
sibling Vevil vehicle, is not publicly confirmed.

\subsection{Wildec Establishment (Liechtenstein): a Belgravia Anstalt}

Wildec Establishment (ROE \texttt{OE004860}) is a Liechtenstein \emph{Anstalt}, a
distinctive establishment entity with no exact UK analogue. Planning records of
the Royal Borough of Kensington and Chelsea name Wildec Establishment as the
client for alterations to 29 Wilton Crescent, a listed six-storey Belgravia house
it purchased in
2018.\footnote{\url{https://docs.planning.org.uk/20211126/115/R2ZG25RPJ1S00/5r6xtq7rp5v29fcb.pdf}}
The property appears in our data at \pounds{}17.5m with a band-D efficiency.

\subsection{AG Bloom LML 2 (Netherlands): an institutional joint-venture vehicle}

AG Bloom LML 2 BV (ROE \texttt{OE002659}, Netherlands, registered 28 October 2022)
shares its naming, registered-office and control-correspondence pattern with the
``AG Bloom'' entities behind the Bloom / Angelo Gordon (now TPG Angelo Gordon)
urban-logistics joint venture announced in December 2021, whose confirmed asset
list includes a Camberwell scheme matching the SE5 address in our
data.\footnote{\url{https://www.thelogisticspoint.com/2021/12/13/bloom-and-angelo-gordon-announce-250mn-ultra-urban-warehouse-joint-venture/}}
The link between this specific vehicle and 61 Lilford Road is inferred from the
shared corporate signatures rather than a single filing; and because the venture
is an ``ultra-urban warehouse'' programme, the residential EPC linkage in our data
(\pounds{}16m, band E) should be read with caution, as it may reflect a mixed-use
or ancillary residential element rather than a purely residential holding.

\subsection{Privilege Oxford (Luxembourg): student accommodation}

Privilege Oxford S.\`a r.l. (ROE \texttt{OE024676}, Luxembourg, registered 8
February 2023) is the registered proprietor of the Slade Park student apartments
in Headington, Oxford (\pounds{}37m in our data), the single largest
foreign-non-profit holding by value. Public reporting confirms an earlier
ownership chain for the same Oxford Brookes scheme (a 2019 sale by Gatehouse Bank
to a Scottish Widows property trust) but does not itself name Privilege Oxford, so
the current title-to-vehicle step is consistent with, but not directly confirmed
by, the public record.

\subsection{Athena Healthcare (Much Hoole) (Isle of Man): a care home}

The proprietor recorded in our data, Athena Healthcare (Much Hoole) Limited (Isle
of Man, UK establishment \texttt{FC036769}, nature of business ``operation of
residential care home''), was renamed Sandstone Care (Much Hoole) Limited on 20
April 2021; the wider Athena Healthcare Group continues as a separate ongoing
operator.\footnote{\url{https://find-and-update.company-information.service.gov.uk/company/FC036769}}
The home, a residential and nursing home near Preston, shows that the
offshore-incorporation universe is not solely prime-residential: operational care
assets are held through the same Crown Dependency structures. (Our register
cross-section predates the rename, hence the Athena name; the site is branded
``Ribble Court'' by the current operator.)

\subsection{Register-only cases}

For four of the shortlisted proprietors only the bare incorporation could be
verified, and no company-level public reporting was located: MRF Managing Trustee
No.\,1 Limited (Jersey, \texttt{OE007919}; the ``Managing Trustee'' form is
consistent with a Jersey Property Unit Trust structure), Balea Stiftung
(Liechtenstein foundation, \texttt{OE017732}), Caerus Limited (Jersey,
\texttt{OE030686}) and Ripley Ventures SA (British Virgin Islands,
\texttt{OE013441}, administered through a BVI registered agent). That the register
yields an address, a jurisdiction and nothing more is itself the point: for a
large share of the offshore universe the public record stops at the incorporation
certificate, which is precisely the transparency gap that a beneficial-ownership
layer (future work) would close.

\section{Caveats and future work}

The ownership cross-section is the January 2025 register, so profiles are
survivorship-conditioned; \texttt{date\_proprietor\_added} reflects the current
owner's registration, not acquisition of the portfolio. Jurisdiction here is
the country of \emph{legal incorporation} of the registered proprietor. A
beneficial-ownership layer (PSC and Register of Overseas Entities linkage, as
built for the London project) would separate offshore incorporation from
offshore control and is left as future work for this branch.

\end{document}
